\documentclass[aspectratio=169,11pt]{beamer}

\usepackage[utf8]{inputenc}
\usepackage[T1]{fontenc}
\usepackage{lmodern}
\usepackage{booktabs}
\usepackage{listings}
\usepackage{xcolor}

\definecolor{GoBlue}{HTML}{00ADD8}
\definecolor{GoDark}{HTML}{073B4C}
\definecolor{GoNavy}{HTML}{14213D}
\definecolor{GoGreen}{HTML}{2A9D8F}
\definecolor{GoYellow}{HTML}{F4A261}
\definecolor{Terminal}{HTML}{101820}
\definecolor{TerminalText}{HTML}{E8F1F2}
\definecolor{SoftGray}{HTML}{F2F5F7}

\usetheme{Madrid}
\useinnertheme{rounded}
\setbeamertemplate{navigation symbols}{}
\setbeamertemplate{itemize item}{\color{GoBlue}$\blacktriangleright$}
\setbeamertemplate{itemize subitem}{\color{GoGreen}$\bullet$}
\setbeamercolor{structure}{fg=GoDark}
\setbeamercolor{palette primary}{bg=GoDark,fg=white}
\setbeamercolor{palette secondary}{bg=GoBlue,fg=GoNavy}
\setbeamercolor{palette tertiary}{bg=GoNavy,fg=white}
\setbeamercolor{frametitle}{bg=GoDark,fg=white}
\setbeamercolor{block title}{bg=GoBlue!25,fg=GoNavy}
\setbeamercolor{block body}{bg=SoftGray,fg=black}
\setbeamercolor{alerted text}{fg=GoYellow!80!black}

\setbeamertemplate{footline}{%
  \leavevmode%
  \hbox{%
    \begin{beamercolorbox}[wd=.76\paperwidth,ht=2.6ex,dp=1.1ex,leftskip=1em]{author in head/foot}%
      \insertshorttitle
    \end{beamercolorbox}%
    \begin{beamercolorbox}[wd=.24\paperwidth,ht=2.6ex,dp=1.1ex,center]{date in head/foot}%
      \insertframenumber{} / \inserttotalframenumber
    \end{beamercolorbox}%
  }%
}

\lstdefinelanguage{Go}{
  morekeywords={break,case,chan,const,continue,default,defer,else,fallthrough,for,func,go,goto,if,import,interface,map,package,range,return,select,struct,switch,type,var},
  sensitive=true,
  morecomment=[l]{//},
  morecomment=[s]{/*}{*/},
  morestring=[b]"
}

\lstdefinestyle{code}{
  basicstyle=\ttfamily\small,
  backgroundcolor=\color{SoftGray},
  keywordstyle=\color{GoDark}\bfseries,
  commentstyle=\color{GoGreen!70!black},
  stringstyle=\color{GoYellow!80!black},
  showstringspaces=false,
  columns=fullflexible,
  keepspaces=true,
  frame=single,
  rulecolor=\color{GoBlue!60},
  frameround=tttt,
  xleftmargin=0.5em,
  xrightmargin=0.5em
}

\lstdefinestyle{terminal}{
  basicstyle=\ttfamily\scriptsize\color{TerminalText},
  backgroundcolor=\color{Terminal},
  showstringspaces=false,
  columns=fullflexible,
  keepspaces=true,
  frame=single,
  rulecolor=\color{GoBlue},
  frameround=tttt,
  xleftmargin=0.5em,
  xrightmargin=0.5em
}

\title[Go: concorrência na prática]{Go: Concorrência Baseada em Goroutines e Channels}
\subtitle{Pesquisa aplicada, exemplos práticos e comparação com C e Python}
\institute[UFCA]{%
  Universidade Federal do Cariri -- UFCA\\
  Centro de Ciência e Tecnologia -- CCT\\
  Bacharelado em Engenharia de Software\\[0.4em]
  Paradigmas de Programação\\
  Prof. Rafael Will Macedo de Araújo
}
\date{}

\begin{document}

\begin{frame}[plain]
  \titlepage
  \begin{center}
    \color{GoDark}\small
    Elder Rayan \textbullet{} Espedio Ramom \textbullet{} Manoel Junio \textbullet{}
    Pedro Yan \textbullet{} Sabrina Alencar \textbullet{} Samuel Wagner \textbullet{} Sebastião Sousa
  \end{center}
\end{frame}

\begin{frame}{Por que Go surgiu?}
  \begin{columns}[T]
    \column{0.57\textwidth}
      \begin{block}{Problema de engenharia}
        Grandes sistemas exigiam builds rápidos, código legível e coordenação eficiente de muitas tarefas de rede em arquiteturas multicore.
      \end{block}
      \begin{itemize}
        \item Projeto iniciado no Google em 2007;
        \item Robert Griesemer, Rob Pike e Ken Thompson;
        \item Go 1.0 lançado em 2012;
        \item simplicidade, tipagem estática e compilação nativa.
      \end{itemize}
    \column{0.39\textwidth}
      \begin{center}
        \setlength{\fboxsep}{10pt}
        \colorbox{GoBlue!15}{\parbox{0.82\linewidth}{\centering
          \textbf{Pergunta central}\\[0.5em]
          Como simplificar software concorrente sem esconder seus riscos?
        }}
      \end{center}
  \end{columns}
\end{frame}

\begin{frame}[fragile]{Goroutines: concorrência leve}
  \begin{columns}[T]
    \column{0.44\textwidth}
\begin{lstlisting}[style=code,language=Go]
var wg sync.WaitGroup
wg.Add(1)

go func() {
    defer wg.Done()
    executarTarefa()
}()

wg.Wait()
\end{lstlisting}
    \column{0.52\textwidth}
      \begin{itemize}
        \item Uma goroutine é uma função executada concorrentemente.
        \item O runtime multiplexa goroutines sobre threads do sistema operacional.
        \item \texttt{WaitGroup} aguarda a conclusão de um conjunto de tarefas.
      \end{itemize}
      \begin{alertblock}{Concorrência não é paralelismo}
        Concorrência estrutura tarefas independentes; paralelismo exige execução simultânea em mais de um núcleo.
      \end{alertblock}
  \end{columns}
\end{frame}

\begin{frame}[fragile]{Channels e interoperabilidade}
  \begin{columns}[T]
    \column{0.47\textwidth}
\begin{lstlisting}[style=code,language=Go]
jobs := make(chan Job)

go func() {
    jobs <- tarefa
    close(jobs)
}()

for job := range jobs {
    processar(job)
}
\end{lstlisting}
    \column{0.49\textwidth}
      \begin{block}{Comunicação explícita}
        \centering Produtor $\longrightarrow$ channel tipado $\longrightarrow$ consumidor
      \end{block}
      \begin{itemize}
        \item envio, recebimento e sincronização;
        \item fechamento indica ausência de novos valores;
        \item opção 2: Go inicia Python e captura sua saída;
        \item uso incorreto ainda pode causar deadlocks e leaks.
      \end{itemize}
  \end{columns}
\end{frame}

\begin{frame}{O mesmo worker pool em três linguagens}
  \small
  \begin{tabular}{@{}p{2.3cm}p{2.7cm}p{3.0cm}p{3.0cm}@{}}
    \toprule
    \textbf{Aspecto} & \textbf{Go} & \textbf{C / pthreads} & \textbf{Python} \\
    \midrule
    Unidade & goroutine & thread POSIX & \texttt{Thread} \\
    Distribuição & channel & fila circular & \texttt{Queue} \\
    Sincronização & channel + \texttt{WaitGroup} & mutex + condições & \texttt{Queue} + \texttt{join} \\
    Gerência & runtime & aplicação + SO & interpretador + SO \\
    CPU paralelo & sim & sim & limitado pelo GIL padrão \\
    Manutenção & coordenação concisa & mais invariantes manuais & API simples \\
    \bottomrule
  \end{tabular}
  \begin{alertblock}{Leitura correta}
    A visualização demonstra coordenação, não desempenho. Benchmark exige ambiente controlado e medições próprias.
  \end{alertblock}
\end{frame}

\begin{frame}[fragile]{Exemplos práticos e pacotes}
  \begin{columns}[T]
    \column{0.47\textwidth}
\begin{lstlisting}[style=terminal]
internal/
|-- exemplos/
|   |-- waitgroup.go
|   `-- channels.go
|-- interop/
|-- estudocaso/
|-- comparativo/
`-- terminal/
\end{lstlisting}
    \column{0.49\textwidth}
      \begin{itemize}
        \item exemplos pequenos isolam cada conceito;
        \item pacotes internos são integrados pelo menu;
        \item produtor--consumidor torna a comunicação observável;
        \item separação favorece leitura, teste e manutenção.
      \end{itemize}
      \begin{block}{Opção 1}
        Demonstra \texttt{WaitGroup} e channel com resultados visíveis no terminal.
      \end{block}
  \end{columns}
\end{frame}

\begin{frame}[fragile]{Estudo de caso: monitor de uptime}
  \begin{columns}[T]
    \column{0.48\textwidth}
      \begin{block}{Fluxo real}
        \centering URLs $\rightarrow$ goroutines HTTP $\rightarrow$ channel $\rightarrow$ terminal
      \end{block}
      \begin{itemize}
        \item uma requisição concorrente por URL;
        \item timeout de 5 segundos;
        \item erros não interrompem as demais tarefas;
        \item channel fechado após todos os workers.
      \end{itemize}
    \column{0.48\textwidth}
\begin{lstlisting}[style=terminal]
ESTUDO DE CASO - MONITOR

[OK] go.dev          HTTP 200
[OK] github.com      HTTP 200
[OK] cloudflare.com  HTTP 200
[ERRO] invalid-url   indisponivel

[OK] Monitoramento concluido
\end{lstlisting}
  \end{columns}
\end{frame}

\begin{frame}[fragile]{CLI: demonstração visual e verificação}
\begin{lstlisting}[style=terminal]
GO - CONCORRENCIA NA PRATICA

  1. Exemplos Basicos
  2. Interoperabilidade Python
  3. Estudo de Caso (Monitor de Uptime)
  4. Comparativo Visual e Verificacao (Go x C x Python)
  0. Sair

> Escolha uma opcao: 4

Go x C:      saidas iguais.
Go x Python: saidas iguais.
[OK] Comparativo completo aprovado: 17984 primos.
\end{lstlisting}
  \begin{center}
    \small CLI $\rightarrow$ \texttt{demonstrar.sh --auto} $\rightarrow$ \texttt{verificar.sh} $\rightarrow$ \texttt{diff}
  \end{center}
\end{frame}

\begin{frame}[fragile]{Reprodutibilidade e documentação}
  \begin{columns}[T]
    \column{0.49\textwidth}
\begin{lstlisting}[style=terminal]
$ go test ./...
$ go test -race ./...
$ go vet ./...
$ ./examples/comparativo/verificar.sh

Go x C:      saidas iguais.
Go x Python: saidas iguais.
\end{lstlisting}
    \column{0.47\textwidth}
      \begin{itemize}
        \item documentação técnica navegável em Markdown;
        \item instruções completas de instalação e execução;
        \item referências oficiais e acadêmicas;
        \item testes Go e Python;
        \item C compilado com avisos tratados como erros.
      \end{itemize}
      \begin{alertblock}{Limite do race detector}
        Ele observa os caminhos executados; não prova ausência de toda data race possível.
      \end{alertblock}
  \end{columns}
\end{frame}

\begin{frame}{Avaliação crítica e aplicações}
  \begin{columns}[T]
    \column{0.48\textwidth}
      \begin{block}{Onde Go se destaca}
        \begin{itemize}
          \item APIs e microsserviços;
          \item pipelines e worker pools;
          \item ferramentas de infraestrutura;
          \item atividades intensivas em rede.
        \end{itemize}
      \end{block}
    \column{0.48\textwidth}
      \begin{alertblock}{O que Go não elimina}
        \begin{itemize}
          \item data races e deadlocks;
          \item vazamentos de goroutines;
          \item cancelamento mal projetado;
          \item necessidade de medir desempenho.
        \end{itemize}
      \end{alertblock}
  \end{columns}
  \begin{center}
    \textbf{A escolha depende do domínio, dos requisitos e da experiência da equipe.}
  \end{center}
\end{frame}

\begin{frame}{Conclusão}
  \begin{center}
    \Large
    Go não elimina a complexidade da concorrência.\\[0.7em]
    \color{GoDark}\textbf{Ele oferece primitivas concisas e integradas que tornam padrões concorrentes mais claros e produtivos.}\\[1em]
    \normalsize\color{black}
    Três linguagens, oito faixas, quatro workers e o mesmo resultado:\\
    \alert{17.984 números primos}.\\[1.2em]
    \textbf{Obrigado! Perguntas?}
  \end{center}
\end{frame}

\begin{frame}{Referências essenciais}
  \footnotesize
  \begin{thebibliography}{9}
    \bibitem{gospec} GO AUTHORS. \emph{The Go Programming Language Specification}. \texttt{https://go.dev/ref/spec}.
    \bibitem{gomemory} GO AUTHORS. \emph{The Go Memory Model}. \texttt{https://go.dev/ref/mem}.
    \bibitem{effectivego} GO AUTHORS. \emph{Effective Go: Concurrency}. \texttt{https://go.dev/doc/effective\_go\#concurrency}.
    \bibitem{hoare} HOARE, C. A. R. Communicating Sequential Processes. \emph{Communications of the ACM}, 1978.
    \bibitem{posix} THE OPEN GROUP; IEEE. \emph{POSIX Threads}.
    \bibitem{python} PYTHON SOFTWARE FOUNDATION. \emph{threading, queue e Global Interpreter Lock}.
  \end{thebibliography}
  \begin{block}{Consulta completa}
    Referências e links de acesso estão em \texttt{docs/referencias.md}.
  \end{block}
\end{frame}

\end{document}
